\section{Conclusion and Future Work}
\label{sec:conclusion}

In this paper, we presented a critical, rigorous sanity check and representational analysis of the widely adopted layer-wise model merging paradigm. By designing and evaluating three distinct diagnostic treatments---intra-task layer shuffling, task-wise spatial averaging, and relative noise injection---we stress-tested the foundational assumption that localized, layer-specific merging coefficients are critical to resolve multi-task weight interference. 

Our multi-seed, dual-optimizer investigation across MNIST, FashionMNIST, CIFAR-10, and SVHN exposes a profound and nuanced dialectical interaction, which we term the \textbf{Overfitting-Optimizer Paradox}. We show that unconstrained test-time adaptation of layer-wise model merging coefficients (52 parameters) on a tiny calibration split of 256 images is highly prone to transductive overfitting. Under a zero-order Adaptive 1+1 Evolution Strategy (1+1 ES), the overfitting takes the form of high-frequency optimization noise that is easily smoothed out by Spatial Averaging (reducing parameters by 92.3\%), which acts as a spatial regularizer and improves average accuracy from $85.07 \pm 0.47\%$ to $85.21 \pm 0.11\%$. Under first-order Adam GD, the gradients allow the optimizer to find a highly precise, delicate configuration of parameters that overfits the calibration statistics. Shuffling this delicate configuration collapses average performance by $5.43\%$ (collapsing CIFAR-10 by $15.69\%$), and spatial averaging collapses CIFAR-10 by $10.35\%$, creating an illusion of functional layer-specificity. However, this optimized model fails to outperform the unoptimized Task Arithmetic baseline ($84.44 \pm 0.37\%$) on the unseen test set, while introducing 4x greater instability and variance, confirming that its layer-specificity is a delicate transductive overfitting artifact rather than a generalizable representation. Furthermore, our representational similarity analysis via linear CKA reveals a key discrepancy, showing that while spatial averaging improves activation similarity on average, activation-level alignment can be highly misleading and does not predict downstream classification performance due to a decoupling between high-level activation correlation and fine-grained decision boundary integrity. Additionally, our noise sensitivity sweep reveals that both optimization regimes are extremely robust, tolerating up to 50\% relative noise on coefficients with negligible performance degradation, confirming that the merging optimization landscape resides in an exceptionally flat basin.

\textbf{Methodological Recommendations:} We offer several guidelines for future model merging and test-time adaptation research:
\begin{enumerate}
    \item \textbf{Mandatory Simple Baselines:} All future publications proposing layer-wise, block-wise, or weight-wise merging parameters must be benchmarked against a properly tuned, single task-wise scalar baseline. Simply comparing against an unoptimized baseline (e.g., Task Arithmetic with $\lambda=0.3$) is methodologically invalid.
    \item \textbf{Diagnostic Control Checks:} Researchers should run shuffling and spatial averaging controls on their learned parameters to verify whether their claimed localized importances are functionally real or merely optimization noise.
    \item \textbf{Representational Verification:} Future works should go beyond validation accuracy and evaluate activation-level similarity to ensure that weight-space updates do not disrupt representational integrity.
\end{enumerate}

\textbf{Limitations and Future Work:} While our results are highly compelling, we acknowledge several critical limitations that restrict the immediate generalizability of our findings. 

First, our evaluation is conducted on saturated, low-resolution vision classification datasets (MNIST, FashionMNIST, CIFAR-10, SVHN) where task vectors reside relatively close to their shared pre-trained CLIP ViT-B/32 initialization. In larger-scale networks, such as 7B+ parameter modern autoregressive decoder-only Large Language Models (LLMs), or on highly complex, high-resolution downstream domains (such as instruction-tuning, cross-modal tasks, or specialized medical and satellite imaging), representational hierarchies are highly distinct. In such scenarios, the task vectors are likely to reside much further from the pre-trained weights in the parameter manifold. Under these conditions of high task divergence, layer-specific weight coordinate adjustments may possess genuine, physical importance to resolve structural and semantic conflicts, and the extreme landscape flatness we observed might no longer hold.

Second, the structural architecture of modern autoregressive decoder LLMs introduces specific layer-wise and block-wise specialization. For instance, representation learning literature has shown that early layers in LLMs primarily capture syntactic and surface-level textual patterns, middle layers encode factual and syntactic-semantic links, and later layers specialize in generation-focused semantic representations. Furthermore, within each transformer block, the attention mechanism (specifically the query, key, and value projection matrices vs. the out-projection matrix) and the feed-forward MLP blocks exhibit highly distinct gradient trajectories and weight dynamics. Merging two highly divergent LLM experts (e.g., a coding expert and a medical expert) might lead to severe layer-specific conflicts (such as middle-layer factual interference or MLP gating conflicts) where simple spatial averaging across all layers collapses the model's generation capability. Exploring the boundaries of the "layer-specificity illusion" across these specific structural sub-components in decoder-only architectures represents a critical, high-impact avenue for future work.

Finally, while we compared zero-order 1+1 ES with first-order Adam GD and evaluated proximity regularization against standard weight decay, investigating the sensitivity of the transductive overfitting threshold to other optimization hyperparameters---such as learning rate schedules, Adam's $\beta_1, \beta_2$ parameters, or standard optimizer weight decay over multiple seeds---remains an open and important direction to help NLP and vision researchers design more robust test-time adaptive merging frameworks.

We hope this study serves as a critical course-correction, steering the model merging community toward a more rigorous, optimizer-aware, and complexity-aware understanding of layer-wise model merging, where localized parameterizations are carefully validated using diagnostic controls.
